\documentclass[UTF8,11pt,a4paper]{ctexart}

\usepackage{amsmath,amssymb,amsthm,mathtools,bm}
\usepackage{booktabs}
\usepackage{enumitem}
\usepackage[margin=2.4cm]{geometry}
\usepackage[colorlinks=true,linkcolor=blue,urlcolor=blue]{hyperref}

\newtheorem{theorem}{定理}[section]
\newtheorem{proposition}[theorem]{命题}
\newtheorem{lemma}[theorem]{引理}
\newtheorem{corollary}[theorem]{推论}
\theoremstyle{definition}
\newtheorem{definition}[theorem]{定义}
\theoremstyle{remark}
\newtheorem{remark}[theorem]{注记}

\newcommand{\R}{\mathbb{R}}
\newcommand{\T}{\mathsf{T}}
\newcommand{\one}{\bm 1}
\newcommand{\dd}{\mathrm d}
\newcommand{\V}{\mathcal V}
\newcommand{\B}{\mathcal B}
\newcommand{\wt}{\widetilde}

\title{Minimum-Balance EC--SAV：面向 Swing 暂态动力学的最小扰动物理能量平衡校正}
\author{}
\date{\today}

\begin{document}
\maketitle

\begin{abstract}
本文记录当前 EC--SAV 新方案的主线思想。与旧版本“只缩放相对速度”或
“kinetic/potential 两分支切换”不同，新方案直接把校正写成一个受
center-of-inertia (COI) 约束的最小扰动物理能量平衡投影。第一步使用线性
SAV--CN predictor；第二步把 predictor 的辅助变量失配
\(
\varepsilon=\V(\wt\delta)+C-\wt r^2
\)
识别为离散物理能量平衡的全部标量缺陷；第三步在步长相关度量
\[
 \|\Delta x\|_h^2
 =h^{-2}\Delta\delta^\T M\Delta\delta
 +\Delta\omega^\T M\Delta\omega
\]
下，沿平衡泛函的 Riesz 梯度做一个一维非线性投影。该方向在一阶意义下
给出满足能量平衡所需的最小状态改动，并自动在正常点偏向速度校正、在
turning point 偏向势能方向，无需人为分支。本文还给出 near-equilibrium
振幅一致估计、turning-point 超一致性、COI 约束和故障切换能量跳变的分析，
并列出配套 Python 仿真的验证指标。本文是当前思路说明和理论草稿，不把尚未
完成的全局误差证明或方法新颖性检索写成既成结论。
\end{abstract}

\section{模型与真实物理能量}

考虑多机 Swing 系统
\begin{equation}
 M\dot\omega+D\omega+\nabla \V(\delta)=0,
 \qquad
 \dot\delta=\omega,
 \label{eq:swing}
\end{equation}
其中 \(M=M^\T\succ0\)，\(D=D^\T\succeq0\)。对经典无损网络，可将
机械功率项并入总势能
\begin{equation}
 \V(\delta)
 =
 U_{\rm net}(\delta)-P_m^\T\delta .
\end{equation}
真实机电能量为
\begin{equation}
 H(\delta,\omega)
 =
 \frac12\omega^\T M\omega+\V(\delta).
 \label{eq:physical-energy}
\end{equation}
连续解满足
\begin{equation}
 \frac{\dd H}{\dd t}
 =
 -\omega^\T D\omega\le0.
 \label{eq:continuous-balance}
\end{equation}

若模型关于统一角度平移不变，且 \(\one^\T P_m=0\)，则
\[
 \V(\delta+c\one)=\V(\delta),
 \qquad
 \one^\T\nabla\V(\delta)=0.
\]
无阻尼时还存在总动量守恒
\[
 \frac{\dd}{\dd t}\bigl(\one^\T M\omega\bigr)=0.
\]

\section{SAV predictor 与标量一致性缺陷}

在所研究的暂态轨道邻域内取 \(C>0\) 使
\(
 \V(\delta)+C>0
\)，定义
\begin{equation}
 q(\delta)=\sqrt{\V(\delta)+C},
 \qquad
 r=q(\delta).
\end{equation}
令
\[
 b(\delta)=\frac{\nabla\V(\delta)}{q(\delta)}=2\nabla q(\delta).
\]
则连续系统等价写成
\begin{align}
 \dot\delta&=\omega,\\
 M\dot\omega+D\omega+r\,b(\delta)&=0,\\
 \dot r&=\frac12 b(\delta)^\T\dot\delta.
\end{align}

设 \(\delta^{*,n+1/2}\) 是二阶显式中点近似，例如
\[
 \delta^{*,n+1/2}
 =
 \frac32\delta^n-\frac12\delta^{n-1},
\]
首步可用
\(
 \delta^n+\frac h2\omega^n
\)。
记
\[
 b^*=
 \frac{\nabla\V(\delta^{*,n+1/2})}
 {\sqrt{\V(\delta^{*,n+1/2})+C}}.
\]
SAV--CN predictor 定义为
\begin{align}
 \frac{\wt\delta^{n+1}-\delta^n}{h}
 &=
 \bar\omega^{n+1/2},
 \label{eq:sav-delta}\\
 M\frac{\wt\omega^{n+1}-\omega^n}{h}
 +D\bar\omega^{n+1/2}
 +\bar r^{n+1/2}b^*
 &=0,
 \label{eq:sav-omega}\\
 \frac{\wt r^{n+1}-r^n}{h}
 &=
 \frac12(b^*)^\T\bar\omega^{n+1/2},
 \label{eq:sav-r}
\end{align}
其中
\[
 \bar\omega^{n+1/2}
 =\frac{\omega^n+\wt\omega^{n+1}}2,
 \qquad
 \bar r^{n+1/2}
 =\frac{r^n+\wt r^{n+1}}2.
\]

\begin{proposition}[SAV predictor 的修改能量律]
若 \(r^n=\sqrt{\V(\delta^n)+C}\)，则 predictor 满足
\begin{equation}
 \frac12(\wt\omega^{n+1})^\T M\wt\omega^{n+1}
 +(\wt r^{n+1})^2-C
 -H^n
 =
 -h(\bar\omega^{n+1/2})^\T D\bar\omega^{n+1/2}.
 \label{eq:sav-modified-law}
\end{equation}
\end{proposition}

定义 SAV consistency defect
\begin{equation}
 \boxed{
 \varepsilon^{n+1}
 =
 \V(\wt\delta^{n+1})+C-(\wt r^{n+1})^2 .}
 \label{eq:epsilon}
\end{equation}
它测量辅助变量与真实势能之间的离散失配。

\section{直接投影离散物理平衡，而不是只投影能量}

对任意候选状态 \((\delta,\omega)\)，定义一步物理平衡泛函
\begin{equation}
 \boxed{
 \B_n(\delta,\omega)
 =
 H(\delta,\omega)-H^n
 +
 h\left(
 \frac{\omega^n+\omega}{2}
 \right)^\T
 D
 \left(
 \frac{\omega^n+\omega}{2}
 \right).}
 \label{eq:balance-functional}
\end{equation}
最终离散解要求
\begin{equation}
 \B_n(\delta^{n+1},\omega^{n+1})=0.
 \label{eq:balance-target}
\end{equation}
于是自动得到
\begin{equation}
 \boxed{
 H^{n+1}-H^n
 =
 -h
 \left(
 \frac{\omega^n+\omega^{n+1}}2
 \right)^\T
 D
 \left(
 \frac{\omega^n+\omega^{n+1}}2
 \right)\le0.}
 \label{eq:exact-discrete-balance}
\end{equation}

\begin{proposition}[缺陷恒等式]
在 predictor 前一步已经 reset \(r^n=q(\delta^n)\) 的条件下，
\begin{equation}
 \boxed{
 \B_n(\wt\delta^{n+1},\wt\omega^{n+1})
 =
 \varepsilon^{n+1}.}
 \label{eq:defect-identity}
\end{equation}
\end{proposition}

\begin{proof}
用真实势能项
\(
 \V(\wt\delta)=\wt r^2-C+\varepsilon
\)
替换式~\eqref{eq:sav-modified-law} 即得。
\end{proof}

这意味着 SAV predictor 对真实离散物理平衡的全部破坏恰好压缩为一个标量
\(\varepsilon\)。

\section{步长相关最小扰动度量}

令状态修正为
\[
 \Delta x=
 \begin{pmatrix}
  \Delta\delta\\
  \Delta\omega
 \end{pmatrix}.
\]
由于 \(\dot\delta=\omega\)，一步内角度误差的自然速度尺度是
\(\Delta\delta/h\)。因此采用
\begin{equation}
 \boxed{
 \|\Delta x\|_h^2
 =
 \frac1{h^2}\Delta\delta^\T M\Delta\delta
 +
 \Delta\omega^\T M\Delta\omega .}
 \label{eq:h-metric}
\end{equation}

在 predictor 点定义
\begin{align}
 f&=\nabla\V(\wt\delta^{n+1}),\\
 \bar{\wt\omega}
 &=\frac{\omega^n+\wt\omega^{n+1}}2,\\
 g_\omega
 &=M\wt\omega^{n+1}+hD\bar{\wt\omega}.
 \label{eq:gomega}
\end{align}
则
\[
 \nabla_\delta\B_n=f,
 \qquad
 \nabla_\omega\B_n=g_\omega.
\]

\subsection{COI 约束}

对旋转对称的多机系统，不允许校正任意改变 common mode。定义
\(M\)-正交 COI 投影
\begin{equation}
 \Pi_M z
 =
 z-\one
 \frac{\one^\T Mz}{\one^\T M\one}.
 \label{eq:coi-projection}
\end{equation}
若不存在整体旋转零模（例如 SMIB），则取 \(\Pi_M=I\)。

受 COI 约束的 Riesz 梯度方向为
\begin{equation}
 \boxed{
 p_h=
 \begin{pmatrix}
 p_\delta\\
 p_\omega
 \end{pmatrix}
 =
 \begin{pmatrix}
 h^2\Pi_M M^{-1}f\\
 \Pi_M M^{-1}g_\omega
 \end{pmatrix}.}
 \label{eq:projection-direction}
\end{equation}
定义
\begin{equation}
 \boxed{
 \chi_h
 =
 f^\T p_\delta+g_\omega^\T p_\omega
 =
 \|p_h\|_h^2\ge0.}
 \label{eq:chi}
\end{equation}

无阻尼且在 COI 坐标中，主导结构退化为
\[
 \chi_h
 =
 h^2\nabla\V^\T M^{-1}\nabla\V
 +
 \wt\omega^\T M\wt\omega.
\]
因此 \(\chi_h\) 同时测量一步内势能力和动能方向对能量缺陷的吸收能力。

\section{一阶最小扰动性质}

线性化平衡约束为
\begin{equation}
 f^\T\Delta\delta
 +
 g_\omega^\T\Delta\omega
 =
 -\varepsilon,
 \label{eq:linearized-balance}
\end{equation}
并要求
\(
 \Delta\delta,\Delta\omega
\)
属于 COI 切空间。

\begin{theorem}[线性化的最小扰动校正]
若 \(\chi_h>0\)，则在所有满足
式~\eqref{eq:linearized-balance} 的 COI-compatible 校正中，
\begin{equation}
 \boxed{
 \Delta x_{\min}
 =
 -\frac{\varepsilon}{\chi_h}p_h}
 \label{eq:linear-min-correction}
\end{equation}
唯一最小化度量~\eqref{eq:h-metric}，且
\begin{equation}
 \boxed{
 \|\Delta x_{\min}\|_h
 =
 \frac{|\varepsilon|}{\sqrt{\chi_h}}.}
 \label{eq:min-bound}
\end{equation}
\end{theorem}

\begin{proof}
对任意可行 \(\Delta x\)，Riesz 表示给出
\[
 -\varepsilon
 =
 \langle p_h,\Delta x\rangle_h.
\]
由 Cauchy--Schwarz，
\[
 |\varepsilon|
 \le
 \sqrt{\chi_h}\,\|\Delta x\|_h.
\]
等号当且仅当 \(\Delta x\) 与 \(p_h\) 共线，结合约束即得到
式~\eqref{eq:linear-min-correction}。
\end{proof}

因此，新方案的核心优点不再表述为“把能量误差变为机器精度”，而是：
在一阶意义下，它以最小的状态扰动修复 SAV predictor 的物理平衡缺陷。

\section{非线性 exact balance projection}

不直接采用线性校正，而是沿最小扰动方向做一维精确投影：
\begin{equation}
 x^{n+1}
 =
 \wt x^{n+1}+\lambda p_h.
 \label{eq:nonlinear-projection}
\end{equation}
定义
\begin{equation}
 F(\lambda)
 =
 \B_n(
 \wt\delta+\lambda p_\delta,\,
 \wt\omega+\lambda p_\omega).
 \label{eq:F}
\end{equation}
则
\begin{equation}
 F(0)=\varepsilon,
 \qquad
 F'(0)=\chi_h.
 \label{eq:F0}
\end{equation}
所以靠近零的小根满足
\begin{equation}
 \boxed{
 \lambda
 =
 -\frac{\varepsilon}{\chi_h}
 +O\!\left(
 \frac{\varepsilon^2}{\chi_h^2}
 \right)}
 \label{eq:lambda-asymptotic}
\end{equation}
（精确余项依赖 \(F''\) 的局部界）。

求得 \(\lambda\) 后 reset
\begin{equation}
 r^{n+1}
 =
 \sqrt{\V(\delta^{n+1})+C}.
 \label{eq:r-reset}
\end{equation}

\begin{remark}
数值实现中不应把 \(\lambda\) 本身当作“校正大小”。真正的状态改动是
\[
 \mu:=\|\Delta x\|_h
 =|\lambda|\sqrt{\chi_h}.
\]
当 \(\chi_h\to0\) 时，\(\lambda\) 可以保持 \(O(h^2)\)，但只要
\(\mu\to0\)，状态映射仍然是正则的。
\end{remark}

\section{SAV defect 的结构估计}

令
\[
 q(\delta)=\sqrt{\V(\delta)+C},
 \qquad
 d=\wt\delta^{n+1}-\delta^n,
 \qquad
 m=\frac{\delta^n+\wt\delta^{n+1}}2.
\]
由 SAV 更新，
\[
 \wt r-r^n
 =
 \nabla q(\delta^*)^\T d.
\]
而中点 Taylor 公式给出
\[
 q(\delta^n+d)-q(\delta^n)
 =
 \nabla q(m)^\T d+O(\|d\|^3).
\]
于是
\begin{equation}
 |q(\wt\delta)-\wt r|
 \le
 C\left(
 \|\delta^*-m\|\,\|d\|
 +\|d\|^3
 \right).
 \label{eq:q-defect-bound}
\end{equation}
在 \(q,\wt r\) 有界且远离零的区域，
\begin{equation}
 \boxed{
 |\varepsilon|
 \le
 C\left(
 \|\delta^*-m\|\,\|d\|
 +\|d\|^3
 \right).}
 \label{eq:epsilon-structural}
\end{equation}

对二阶 predictor，
\[
 \|\delta^*-m\|=O(h^2),
 \qquad
 \|d\|=O(h),
\]
故一般点
\begin{equation}
 \varepsilon=O(h^3).
 \label{eq:regular-epsilon}
\end{equation}
在非退化 turning point，\(\omega(t_n)=0\) 而加速度非零，此时
\(
 d=O(h^2)
\)，因而
\begin{equation}
 \boxed{\varepsilon=O(h^4).}
 \label{eq:turning-epsilon}
\end{equation}
这个额外一阶来自状态位移本身的缩小，不是某个单机算例中的偶然抵消。

\section{稳定平衡点附近的振幅一致估计}

在 COI reduced space 中设
\[
 \nabla\V(\delta_e)=0,
 \qquad
 K=\nabla^2\V(\delta_e)\succ0.
\]
令
\[
 y=\delta-\delta_e,\qquad v=\omega,
\]
并定义小扰动振幅
\begin{equation}
 \rho^2=\|v\|_M^2+\|y\|_K^2.
 \label{eq:rho}
\end{equation}
同时定义步长相关动力学尺度
\begin{equation}
 \boxed{
 \sigma_h^2
 =
 \|v\|_M^2
 +
 h^2
 \|\nabla\V(\delta)\|_{M^{-1}}^2.}
 \label{eq:sigma}
\end{equation}
在稳定平衡点小邻域，
\[
 c h\rho\le\sigma_h\le C\rho.
\]

由于
\[
 \V(\delta_e+y)
 =
 \V(\delta_e)+\frac12y^\T Ky+O(\|y\|^3),
\]
有
\[
 q(\delta_e+y)
 =
 q_e+\frac{1}{4q_e}y^\T Ky+O(\|y\|^3),
\]
故 \(\nabla q(\delta_e)=0\)。结合
式~\eqref{eq:q-defect-bound} 以及二阶 predictor 的
\[
 \|d\|\le Ch\sigma_h,
 \qquad
 \|\delta^*-m\|\le Ch^2\rho,
\]
得到工作估计
\begin{equation}
 \boxed{
 |\varepsilon|
 \le
 C h^3\rho\sigma_h.}
 \label{eq:near-equilibrium-epsilon}
\end{equation}
因此
\[
 \varepsilon=O(\rho^2h^3).
\]
在 turning point，\(\sigma_h=O(h\rho)\)，进一步有
\begin{equation}
 \boxed{
 \varepsilon=O(\rho^2h^4).}
 \label{eq:near-equilibrium-turning}
\end{equation}

另一方面，在充分小步长下
\begin{equation}
 c\sigma_h^2\le\chi_h\le C\sigma_h^2.
 \label{eq:chi-sigma}
\end{equation}
于是
\begin{equation}
 \frac{|\varepsilon|}{\sqrt{\chi_h}}
 \le Ch^3\rho,
 \qquad
 \frac{|\varepsilon|}{\chi_h}
 \le Ch^2.
 \label{eq:ratio-bounds}
\end{equation}

\begin{proposition}[平衡点附近的可去退化]
在上述局部假设以及标量小根存在的条件下，
\begin{equation}
 \boxed{
 \|x^{n+1}-\wt x^{n+1}\|_h
 \le Ch^3\rho.}
 \label{eq:near-equilibrium-correction}
\end{equation}
因此当 \(\rho\to0\) 时，状态校正连续趋于零。即使
\(\lambda\) 本身在 turning-point 方向仅为 \(O(h^2)\)，最终状态映射仍无
振幅奇异。
\end{proposition}

精确平衡点处
\(
 \rho=\varepsilon=\chi_h=0
\)，算法直接定义 correction 为零。

\section{二阶精度与线性模态的当前结论}

若 SAV predictor 的局部截断误差是 \(O(h^3)\)，且上述 projection 满足
\[
 \|\Delta x\|_h=O(h^3),
\]
则校正不降低 predictor 的二阶全局精度。更具体地，在小扰动线性化
\[
 M\ddot y+D\dot y+Ky=0
\]
附近，若 SAV predictor 的一步矩阵为 \(A_h\)，则校正后的一步映射满足
\begin{equation}
 \boxed{
 A_h^{\rm MB-EC-SAV}
 =
 A_h^{\rm SAV}+O(h^3).}
 \label{eq:linear-map}
\end{equation}
因此方法保持二阶一致性，但可能改变 \(O(h^3)\) 的 phase-lag 和 damping
误差系数。这个 \(O(h^3)\) 系数正是下一阶段需要做单模态放大因子分析的对象，
因为“能量闭合”本身并不保证轨迹相位更准确。

\section{故障切换与真实能量 jump}

暂态稳定中网络参数在 fault-on、clearing、post-fault 三个阶段分段变化。
设切换前后势能为 \(\V^-\) 和 \(\V^+\)。理想瞬时开关下
\[
 \delta(t_s^+)=\delta(t_s^-),
 \qquad
 \omega(t_s^+)=\omega(t_s^-),
\]
但物理能量应真实跳变
\begin{equation}
 \boxed{
 H^+-H^-
 =
 \V^+(\delta_s)-\V^-(\delta_s).}
 \label{eq:switch-jump}
\end{equation}
算法不应把这个 jump 当作数值漂移去校正。正确处理是保持
\((\delta,\omega)\) 连续，并按新网络重新初始化
\[
 r^+
 =
 \sqrt{\V^+(\delta_s)+C}.
\]
同时在 event 后重启二阶 extrapolation，避免用跨拓扑的旧 predictor history。

若每一连续阶段都满足式~\eqref{eq:exact-discrete-balance}，则离散全局能量账本为
\begin{equation}
 H^N-H^0
 =
 -\sum_n
 h_n(\bar\omega_n)^\T D_n\bar\omega_n
 +
 \sum_k
 \left[
 \V_k^+(\delta_k)-\V_k^-(\delta_k)
 \right].
 \label{eq:piecewise-energy}
\end{equation}

\section{算法摘要}

每一步执行：
\begin{enumerate}[label=\arabic*.]
 \item 由式~\eqref{eq:sav-delta}--\eqref{eq:sav-r} 做一次线性 SAV predictor；
 \item 计算
 \[
  \varepsilon
  =
  \V(\wt\delta)+C-\wt r^2;
 \]
 \item 计算 \(f,g_\omega,p_h,\chi_h\)；
 \item 若 \(|\varepsilon|\) 已在机器误差范围内，直接接受 predictor；
 \item 否则以
 \[
  \lambda_0=-\varepsilon/\chi_h
 \]
 为初值，解唯一一维方程
 \[
  F(\lambda)=0;
 \]
 \item 更新
 \[
  (\delta^{n+1},\omega^{n+1})
  =
  (\wt\delta,\wt\omega)+\lambda p_h;
 \]
 \item reset
 \[
  r^{n+1}=\sqrt{\V(\delta^{n+1})+C}.
 \]
\end{enumerate}

整个 correction 不解高维非线性系统；非线性部分只有一个标量 root。

\section{配套 Python 仿真应该验证什么}

新脚本 \texttt{minimum\_balance\_ecsav.py} 设计为本地运行，至少输出：

\begin{enumerate}[label=(\arabic*)]
 \item predictor identity
 \[
  \B_n(\wt x)-\varepsilon;
 \]
 \item regular point 的 \(\varepsilon=O(h^3)\)；
 \item turning point 的 \(\varepsilon=O(h^4)\)；
 \item
 \[
  \chi_h,\quad
  \sigma_h,\quad
  \lambda,\quad
  \|\Delta x\|_h;
 \]
 \item exact corrected balance residual
 \[
  |\B_n(x^{n+1})|;
 \]
 \item near-equilibrium 振幅缩放：
 \[
  \varepsilon/(\rho^2h^3),
  \qquad
  \|\Delta x\|_h/(\rho h^3);
 \]
 \item SMIB 全局二阶收敛；
 \item IEEE 9-bus fault/clearing CCT，并与仓库已有 RK4、midpoint、SAV、
 old EC--SAV 结果对比。
\end{enumerate}

真正优先观察的是 state error、CCT error 和 correction norm，而不只是
“能量残差是否达到 \(10^{-15}\)”。

\section{当前理论边界与下一步}

目前最重要的剩余理论问题有：
\begin{enumerate}[label=(\arabic*)]
 \item 给出 \(F''\) 关于 \(\chi_h\) 的统一局部界，从而把小根存在唯一性写成
 完整定理；
 \item 对单模态
 \[
  \ddot y+2\zeta\Omega\dot y+\Omega^2y=0
 \]
 计算校正前后的 amplification factor 到 \(h^3/h^4\)，明确 phase-lag
 系数到底改善还是恶化；
 \item 将局部投影误差估计与 event restart 拼接成含 fault switching 的
 全局二阶误差证明；
 \item 对大扰动、跨越势垒甚至 pole slipping 的情况，处理普通平方根 SAV
 所需的局部下界 \( \V+C>0 \) 问题。
\end{enumerate}

因此当前方案最准确的定位是：\emph{一个以 SAV 线性 predictor 为基础、
直接恢复真实离散耗散平衡、在 COI 约束下具有一阶最小状态扰动性质的
一维 nonlinear balance projection 方法。}

\end{document}
